\chapter{Présentation du projet}

Dans le cadre du module \module, nous nous intéressons à la plateforme \nomProjet\ (\intitule).

Les étudiants ont souvent besoin de livres, de calculatrices, de fournitures scolaires et d'autres équipements pour leurs études. Cependant, l'achat de produits neufs peut représenter une dépense importante pour certains d'entre eux.

Parallèlement, d'autres étudiants possèdent des produits encore utilisables dont ils n'ont plus besoin et qu'ils souhaitent vendre.

Actuellement, les annonces sont dispersées entre les groupes WhatsApp, les réseaux sociaux et le bouche-à-oreille, ce qui rend difficile la recherche de produits, la comparaison des offres et la mise en relation entre acheteurs et vendeurs.

L'objectif est de centraliser ces annonces dans un espace unique afin de faciliter l'achat et la vente de produits de seconde main entre étudiants. Cette idée devra être confirmée auprès de quelques étudiants afin de vérifier qu'elle correspond à une difficulté réellement rencontrée.
